\section{Complete checkers and measured simplicity}
\label{sec:comparison}

\subsection{The main typed checker and matched controls}

The recommended implementation of $\mathcal R$ is
\path{code/orbit_checker.py}. It reads exactly a JSON object with
\texttt{goal} and \texttt{proof} fields, verifies that the goal is a
closed formula, synthesizes the proof's conclusion, and requires exact
agreement. Term and formula tags, arities, integer indices, binder scopes,
context membership, and eigenvariable conditions are checked. A successful
run exits normally and prints exactly \texttt{ACCEPT} followed by a
newline. Any exception or resource failure rejects the input.

The two comparison modules implement the pinned eight-rule $\mathcal N$
and seven-rule $W$ calculi with the same syntax, substitution operations,
validation, and file interface. They are newly written matched controls;
they are not represented as Python files published upstream. All helpers
and the JSON interface are included in each standalone module.

\begin{center}
\small
\begin{tabular}{@{}lrrrr@{}}
\toprule
Complete typed checker & Tokens & AST & Bytes & Gzip 9 \\
\midrule
Adopted $\mathcal N$, matched & 1,235 & 1,407 & 4,692 & 1,199 \\
New orbit $\mathcal R$ & 1,206 & 1,366 & 4,553 & 1,211 \\
Known $W$, matched & 1,173 & 1,334 & 4,456 & 1,198 \\
\bottomrule
\end{tabular}
\end{center}

These numbers use CPython 3.12.14 and zlib 1.3.2. Genuine docstrings are
removed from the parsed AST before \texttt{ast.unparse} normalization;
comments and formatting are consequently absent. Bytes include the final
newline. Gzip uses compression level 9 with \texttt{mtime=0} and includes
its wrapper. The evidence also records raw source bytes, raw gzip sizes,
and both source hashes. Gzip is informational and is not minimized.

Relative to the adopted eight-rule control, the new checker saves
139 normalized bytes (about 2.96\%), 41 AST nodes, and 29 significant
tokens. It nevertheless costs 97 bytes, 32 AST nodes, and 33 tokens more
than the matched $W$ checker. Its normalized gzip is also larger than
both controls. The uniform introduction and direct substitution theorem
therefore do not establish a new minimum under the user's implementation
metric.

\subsection{A separate compact implementation improvement}

The archive also contains three compact implementations of the
\emph{earlier guarded-source calculus}. They are not implementations of
the new orbit calculus. Their comparison contract is the same as that of
the authenticated incoming guarded-source package: canonical guarded
public statements, with broader administrative proof data. The compact
programs do not themselves reject every noncanonical public statement.
The main typed checker, in contrast, validates its complete Eq/$Q$/$D$
public grammar. These are different input contracts.

\begin{center}
\small
\begin{tabular}{@{}lrrrr@{}}
\toprule
Compact complete module & Tokens & AST & Bytes & Gzip 9 \\
\midrule
Previous token profile & 257 & 285 & 568 & 338 \\
New token profile & \textbf{256} & 283 & 593 & 346 \\
Previous AST profile & 261 & 279 & 603 & 345 \\
New AST profile & 260 & \textbf{278} & 602 & 344 \\
Previous byte profile & 264 & 301 & 522 & 328 \\
New byte profile & 263 & 300 & \textbf{521} & 327 \\
\bottomrule
\end{tabular}
\end{center}

Thus the inspected attained coordinate envelope changes from
\[
 (257\text{ tokens},279\text{ AST nodes},522\text{ bytes})
 \quad\hbox{to}\quad
 (256,278,521).
\]
Each triple is attained by three different programs. No single module
attains all three coordinates. The token profile actually gets longer
in bytes and gzip, which is why reporting only the coordinate envelope
would conceal relevant information.

The byte and AST changes replace the initial administrative truth
\texttt{[]} by raw integer \texttt{0}, and change neutral major fields
in emitted certificates accordingly. These are both spellings of a
true administrative resource, not new numerical axioms. At nonneutral
major premises the original rule tests are unchanged. At a neutral
major, both representations select the same computation or introduction
branch and contribute no mathematical assumption. The old canonical
proof-production construction therefore compiles with the new neutral
spelling, while soundness uses the same interpretation of retained
resources. No complete equivalence claim is made about every malformed
or noncanonical raw byte stream.

The token profile instead turns the recursive continuation into a
\texttt{while} loop and uses an immutable tuple context. Its recursive
major-premise call is retained. The loop state is exactly the old
continuation state: the current goal, retained context, and global fresh
index. Induction on consumed packets identifies the two executions.
The global certificate convention is unchanged.

These operations compose techniques already present in the incoming
packages. They do not constitute a new arithmetic basis or a new logical
axiomatization. Their strict implementation gain is real but very small.
The frozen predecessors, their hashes, the changed certificates, and the
entire new programs are supplied for inspection.

\subsection{What the executable evidence checks}

The recorded reproduction run has 270 process expectations: 122 expected
acceptances and 148 expected rejections, each checked in ordinary or
optimized Python as specified in the manifest. Two additional tests
exercise open sequents with displayed assumptions, and four controls
exercise capture-avoiding substitution. There are zero mismatches.

Positive examples cover both computation equations, reflexivity derived
from computation and transport, equation induction whose hypothesis is
actually used, zero- and nonzero-source orbit specializations, compound
case introduction, and closure extraction through an unused witness.
In particular, a complete closed certificate proves
$I(3,0;u.Su)=3$ using the new induction construction.

Negative examples include a false equation, wrong computation endpoints,
an incorrect orbit antecedent, a future witness, an escaping existential
eigenparameter, unbound iterator state, negative hypothesis indices,
noninteger variable indices, attempted compound induction, attempted
primitive compound transport, malformed instructions, and an open public
goal. These address the rules that could otherwise strengthen the theory
or break scoping.

One open test derives $0=1$ from a \emph{false progressive parity
assumption}. It shows that the new closure retains its progressive
meaning; it is not a closed contradiction proof. Another attempts to use
a compound closure progressively and is rejected. The full assumption
lists and their scope are recorded in \path{tests/open_sequents.json}.

Finite checks support the implementation, while the written
metatheorems establish the universal mathematical claim. No number of
these finite controls would by itself prove soundness, completeness,
bi-interpretability, or a minimality theorem.
